\section{Introduction}
Professional work is organized around persistent responsibilities: a researcher collects evidence, an analyst transforms it, a reviewer checks the resulting claims, and a writer produces a deliverable. The organization carries procedures, tools, and handoff conventions from one assignment to the next. For a large language model (LLM) system, reproducing this arrangement requires more than assigning expert names. An analyst needs the correct computation tools and input resources; a reviewer needs the actual revision under review; and improvements discovered during one assignment must survive beyond its conversation.

We introduce \emph{Evolving Expert Organizations} (EEOs), a framework for representing, specializing, executing, and revising reusable agent organizations. Its implementation in \oc{} treats an expert organization as a versioned package containing role instructions, capability allocations, task-level collaboration graphs, and domain resources. A task binds to an exact package revision. An outer improvement process constructs a new revision from recorded experience and compares it with its parent on matched tasks. The resulting object of reuse is the organization, including its procedures and executable resource closure.

Consider a research report whose calculations are correct but whose prose uses an earlier result. Repairing that report addresses one assignment. Refining the writer's procedure so that each numerical claim is explicitly bound to its canonical result and revision-specific review targets a recurring organizational failure. The latter is an organizational update only when it is retained in a package and applied to subsequent tasks. Whether it improves those tasks is a separate, measurable question. This distinction connects the representation, update mechanism, and evaluation throughout the paper.

The paper makes three contributions. First, we define a portable organization representation that separates reusable expertise from task-local state and gives domain specialization an explicit resource and output contract. Second, we define a constrained revision protocol that connects the mutation surface, matched parent--candidate evaluation, and retained version-specific evidence. Third, we report a preliminary AutomationBench execution result and identify the evidence needed to attribute future changes to collaboration, specialization, or retained updating. Our technical analysis concerns the organizational mechanisms; the available experiment characterizes Mission Base and does not isolate those mechanisms.

\section{Related Work}\label{sec:related}
\paragraph{Expert collaboration.}
Role specialization and structured communication are established design choices. AutoGen supplies configurable conversational agents, and MetaGPT uses professional roles and operating procedures to coordinate software work \citep{autogen,metagpt}. EEO makes the portable package the unit of specialization and revision: roles are coupled to their capability references, supporting resources, and handoff declarations. Neither the use of expert roles nor a dependency graph alone constitutes the contribution.

\paragraph{Automatic agent and workflow design.}
EvoAgent mutates agent settings, selects diverse offspring with an LLM quality check, and integrates their task outputs \citep{evoagent}. AFlow represents workflows in code and uses execution feedback in a variant of Monte Carlo tree search \citep{aflow}. ADAS uses a meta agent to program new agents from an archive of prior designs and evaluations, including transfer across tasks, domains, and models \citep{adas}. These results motivate treating organizational revision as search, but they preclude claiming that automatic design or cross-task reuse originates here. EEO studies a constrained package search space that carries domain resources and capability allocations together with instructions. Its current generator edits the procedural layer; broader structural edits are represented and checked by the package integrity interface.

\paragraph{Reflective prompt optimization.}
GEPA (Genetic-Pareto) already uses execution feedback to revise module prompts, retain candidates, and optimize held-out performance with fixed model weights \citep{gepa}. EEO's text configuration therefore does not introduce reflection-based optimization. Its contribution concerns professional organization packaging, explicit resource boundaries, and attributable revision retention. A matched optimizer comparison on this substrate is needed to attribute benefits to the particular revision policy.

\paragraph{Learning from experience.}
SiriuS retains successful reasoning trajectories, augments unsuccessful ones, and performs supervised fine-tuning of individual agent models \citep{sirius}. EEO holds the foundation model fixed and retains improvements in organization packages. The two update different objects and can be combined, but an experience archive by itself does not establish learning: a retained package change and its measured transfer are the relevant evidence here.

\paragraph{Procedural memory and versioned improvement.}
LEGOMem distills successful trajectories into full-task and agent-specific memories and retrieves them for the orchestrator and task agents \citep{legomem}. It directly addresses reusable procedural knowledge in multi-agent workflows. EEO instead revises a package's standing instructions and resources; cross-task procedural reuse is common to both. AgentDevel externalizes agent revision into a versioned development process with executable diagnosis, candidate evaluation, and regression-aware selection \citep{agentdevel}. Thus, external improvement and traceable version retention are also established ideas. EEO's specific design couples the professional organization to an explicit text mutation surface and inherited resource declarations. This combination motivates comparative study; it does not establish superiority over memory retrieval or other revision policies.

\paragraph{Enterprise harness evolution.}
StarHarness revises environment-specific harnesses with fixed model weights, a candidate ledger, development traces, and separate selection and holdout roles; its experiments include AutomationBench Finance \citep{starharness}. Its editable space includes executable interfaces and agent-loop configuration. EEO's present text study holds those resources fixed. A fair comparison must disclose this difference in editable space. Enterprise harness evolution and its evaluation on AutomationBench are prior contributions; the preliminary Mission Base result does not supply such a head-to-head comparison.

\paragraph{Execution substrate.}
Agent environments and interfaces, exemplified by SWE-agent and OpenHands, determine how agents interact with tools and artifacts \citep{sweagent,openhands}. The OpenHands Software Agent SDK also provides persistent event-oriented execution \citep{ohsdk}. EEO uses exact package binding and revision-specific artifacts to keep organizational comparisons attributable. These supporting mechanisms do not by themselves imply better reasoning or semantic correctness.
